\documentclass[12pt]{article}
\usepackage[margin=1in]{geometry}
\usepackage{lmodern}\usepackage[T1]{fontenc}
\usepackage{amsmath,amssymb}
\usepackage{booktabs}
\usepackage{hyperref}
\usepackage{natbib}
\usepackage{microtype}

\title{\textbf{Subdivision-Respecting Redistricting:\\
County-Sticky Edge Weights}\\[0.4em]
\large B.10 --- Algorithm Track}
\author{Gio Della-Libera \\ \textit{Apportionment Research Group}}
\date{May 2026}

\begin{document}
\maketitle

\begin{abstract}
Most state constitutions require congressional maps to ``preserve political
subdivisions'' --- counties and municipalities --- as a secondary criterion
after equal population.  Existing algorithmic approaches either ignore this
requirement or impose hard constraints that degrade compactness.
We introduce \emph{county-sticky} edge weighting: edges between census tracts
within the same county receive a multiplier $\alpha_c \geq 1$, making
intra-county cuts more expensive without prohibiting them.
On all 50 states, $\alpha_c = 3.0$ reduces county splits by $34\%$ relative
to the baseline (B.2) while maintaining $97\%$ of the compactness improvement.
This soft-constraint approach is the weight mode exposed as
\texttt{--weights-override county} in the \texttt{redist} CLI.
\end{abstract}

\section{Introduction}

State constitutional redistricting criteria typically list criteria in priority
order: (1) equal population, (2) VRA compliance, (3) contiguity, (4)
compactness, (5) preservation of political subdivisions.
Items (1)--(4) are addressed by the baseline algorithm (B.2).
Item (5) --- county preservation --- requires explicit treatment.

Two naive approaches fail:
\emph{Hard constraints} (forbidden edges across county boundaries) make large
counties unpartitionable when $k > $ county count, breaking contiguity.
\emph{No treatment} (ignore counties) produces maps that split counties
unnecessarily, violating state constitutional requirements.

We propose a middle path: \emph{county-sticky weights}.

\section{County-Sticky Weighting}

\subsection{Formulation}

Let $c(v)$ denote the county of tract $v$.  For each edge $(u, v) \in E$,
define the county-sticky weight:
\[
  w_E'(u, v) = w_E(u, v) \cdot
  \begin{cases}
    \alpha_c & \text{if } c(u) = c(v) \\
    1 & \text{otherwise}
  \end{cases}
\]
where $w_E(u, v)$ is the shared boundary length (B.2 baseline) and
$\alpha_c \geq 1$ is the county stickiness parameter.

Higher $\alpha_c$ penalises intra-county cuts more, preserving county
integrity at the cost of compactness.

\subsection{Parameter selection}

We sweep $\alpha_c \in \{1.5, 2.0, 3.0, 5.0, 10.0\}$ on all 50 states.
The Pareto frontier of (county splits, compactness) shows a sharp elbow
at $\alpha_c = 3.0$: further increases reduce county splits by $< 5\%$ while
compactness falls by $> 8\%$.

\section{Results}

\begin{table}[h]
\centering
\caption{County preservation vs compactness ($\alpha_c = 3.0$, 2020 Census)}
\begin{tabular}{lrrrr}
\toprule
Metric & Baseline (B.2) & County-Sticky & Change \\
\midrule
Mean Polsby--Popper & 0.361 & 0.350 & $-3.0\%$ \\
County splits & 487 & 323 & $-34\%$ \\
Multi-county districts & 312 & 198 & $-37\%$ \\
Max deviation from ideal pop. & 0.41\% & 0.44\% & $+0.03$pp \\
\bottomrule
\end{tabular}
\end{table}

County preservation improves substantially while compactness falls only 3\%
--- well within the margin of improvement over enacted maps (+22\% vs B.2
baseline).

\subsection{State-level highlights}

Iowa ($k{=}4$, 99 counties): county splits drop from 12 to 3.
The four-district map nearly respects the traditional four-region Iowa
geography (northwest, northeast, southwest, southeast).

Texas ($k{=}38$, 254 counties): county splits drop from 89 to 61.
The 15 largest counties still require splitting ($\geq 4$ districts each),
but medium counties are now largely preserved.

\section{Connection to the Districting Integrity Act}

The DIA's algorithm specification uses county-sticky weights
($\alpha_c = 3.0$) as a default, satisfying state constitutional
subdivision-preservation requirements without hard constraints.
The parameter $\alpha_c$ is published in the TIGER adjacency graph
metadata, making the weighting fully reproducible.

\section{Conclusion}

County-sticky edge weights with $\alpha_c = 3.0$ reduce county splits by 34\%
at a compactness cost of 3\%, satisfying state constitutional
subdivision-preservation criteria without sacrificing the core compactness
gains of edge-weighted bisection.
This approach is implemented as \texttt{--weights-override county} in the
\texttt{redist} CLI.

\bibliographystyle{plainnat}
\bibliography{../references}

\end{document}
